\documentclass[a4paper,twoside]{article}

\usepackage{morewrites}

\usepackage[svgnames,dvipsnames,x11names]{xcolor}
\usepackage{graphicx}
\usepackage[english]{babel}
\usepackage[colorlinks=true, linkcolor=NavyBlue, urlcolor=NavyBlue, citecolor=NavyBlue]{hyperref}
\usepackage[a4paper]{geometry}
\usepackage{ts-fonts}
\usepackage{csquotes}
\usepackage{amsmath}
\usepackage{float}
\usepackage{seqsplit}
\usepackage{ts-callouts}
\usepackage{ts-code}

\usepackage{lastpage}
\usepackage{refcount}
\makeatletter
\def\ps@plain{
  \let\@oddhead\@empty
  \let\@evenhead\@empty
  \def\@oddfoot{
    \hfil
    \ifnum\getpagerefnumber{LastPage}>1\relax
      \thepage
    \fi
    \hfil}
  \let\@evenfoot\@oddfoot
}
\makeatother

\title{References}
\author{}\date{}

\begin{document}

\maketitle

\thispagestyle{plain}
We define several types of references that can be used throughout the documentation:

\begin{description}
\item[{ Internal References }] These are links that point to other sections within the same document or to other documents within the same project.
\item[{ External References }] These are links that point to resources outside of the current project, such as websites or external documents.
\item[{ Bibliographic References }] These are citations that refer to external publications, articles, or books. They are often formatted using a specific citation style (e.g., APA, MLA) and may include a bibliography section at the end of the document.
\item[{ Footnotes }] Footnotes provide additional information or citations without cluttering the main text. They are typically indicated by a superscript number in the text, with the corresponding footnote text provided at the bottom of the page or section.
\item[{ Equations }] Mathematical equations can be included in the document using LaTeX syntax. Equations can be labeled and referenced throughout the text.
\item[{ Tables }] Tables are used to present data in a structured format with rows and columns. They can be labeled and referenced within the document.
\item[{ Figures }] Figures are images, charts, or diagrams included in the document. They can be labeled and referenced throughout the text.
\item[{ Listings/Code Blocks }] Code blocks are used to display code snippets in various programming languages. They can be labeled and referenced within the document.
\item[{ Tags/Index }] Tags or index entries allow you to associate keywords with specific sections or topics in the document, making it easier to locate related information.

\end{description}
\section{Internal References}\label{internal-references}

You can reference another section in the same document or cross-link to other files in the project.

Linking to another file? TeXSmith targets the destination's main heading and drops a proper hyperlink---handy for navigation-friendly PDFs without any manual tinkering.

\begin{code}{markdown}{}{}
\begin{Verbatim}[commandchars=\\\{\},breaklines, breakanywhere, commandchars=\\\{\}]
See the [\PY{n+nt}{Code Examples}](\PY{n+na}{code.md}) for more details.
\end{Verbatim}
\end{code}
Skip the link text and TeXSmith injects the section number for the print build automatically.
\texttt{See the section [](code.md) for more details.
}

\begin{code}{markdown}{}{}
\begin{Verbatim}[commandchars=\\\{\},breaklines, breakanywhere, commandchars=\\\{\}]
\PY{g+gu}{\PYZsh{}\PYZsh{} Section Title \PYZob{}\PYZsh{}sec:section\PYZhy{}title\PYZcb{}}

Placeholder text that other sections can reference.

\PY{g+gu}{\PYZsh{}\PYZsh{} Other Section}

Check section @[sec:section\PYZhy{}title] for more details.
\end{Verbatim}
\end{code}
The brackets are optional when the label is a single word made of letters,
digits and \texttt{\_ : . -\allowbreak{}} characters, so \texttt{@sec:section-\allowbreak{}title} works too. Trailing
sentence punctuation stays out of the label (\texttt{…voir @sec:intro.} references
\texttt{sec:intro}). E-mail addresses and \texttt{@} inside words or URLs are left alone;
write \texttt{\textbackslash{}@} to force a literal \texttt{@} where the shorthand would otherwise apply.

\subsection{Custom counters}\label{custom-counters}

Document-specific series (findings, requirements, bugs) declared under
\texttt{counters:} in the front matter reuse the very same \texttt{@prefix:key} shorthand,
resolving to their formatted number instead of a section number. See
\hyperref[snippet:<HASH>]{Custom counters}.

\subsection{Cross-document references}\label{cross-document-references}

\texttt{@alias:key} resolves against another document's published inventory, declared
under \texttt{crossrefs:} in the front matter. See
\hyperref[snippet:<HASH>]{Cross-document references}.

\subsection{Autorefs}\label{autorefs}

When the \texttt{mkdocs-\allowbreak{}autorefs} extension is enabled, you can use the \texttt{[text][label]} syntax to generate automatic references to headings.

\section{External References}\label{external-references}

Reference external resources (HTTP/HTTPS) with vanilla Markdown link syntax:

\begin{code}{markdown}{}{}
\begin{Verbatim}[commandchars=\\\{\},breaklines, breakanywhere, commandchars=\\\{\}]
For more information, visit the [\PY{n+nt}{TeXSmith Website}](\PY{n+na}{https://texsmith.org}).
You can also check our GitHub repository at https://github.com/yves\PYZhy{}chevallier/texsmith.
\end{Verbatim}
\end{code}
Printed output uses the usual \LaTeX{} link commands:

\begin{code}{latex}{}{}
\begin{Verbatim}[commandchars=\\\{\},breaklines, breakanywhere, commandchars=\\\{\}]
For more information, visit the \PY{k}{\PYZbs{}href}\PY{n+nb}{\PYZob{}}https://texsmith.org\PY{n+nb}{\PYZcb{}}\PY{n+nb}{\PYZob{}}TeXSmith Website\PY{n+nb}{\PYZcb{}}.
You can also check our GitHub repository at \PY{k}{\PYZbs{}url}\PY{n+nb}{\PYZob{}}https://github.com/yves\PYZhy{}chevallier/texsmith\PY{n+nb}{\PYZcb{}}.
\end{Verbatim}
\end{code}
\section{Bibliographic References}\label{bibliographic-references}

Markdown lacks native bibliography support, so TeXSmith reuses the footnote syntax and BibTeX/front matter keys. See the documentation on \hyperref[snippet:<HASH>]{Bibliography management} for more details.

\begin{code}{markdown}{}{}
\begin{Verbatim}[commandchars=\\\{\},breaklines, breakanywhere, commandchars=\\\{\}]
\PYZhy{}\PYZhy{}\PYZhy{}
bibliography:
\PY{g+gu}{  einstein1905: https://doi.org/10.1002/andp.19053221004}
\PY{g+gu}{\PYZhy{}\PYZhy{}\PYZhy{}}
Einstein\PYZsq{}s theory of relativity revolutionized physics. [\PYZca{}einstein1905]
\end{Verbatim}
\end{code}
\section{Footnotes}\label{footnotes}

Use footnotes to park side comments without cluttering the main text. Markdown marks them with superscript numbers; the rendered document moves the details to the bottom of the page or section.

\begin{code}{markdown}{}{}
\begin{Verbatim}[commandchars=\\\{\},breaklines, breakanywhere, commandchars=\\\{\}]
This is a sample sentence with a footnote.[\PYZca{}1]

[\PY{n+nl}{\PYZca{}1}]: \PY{n+na}{This is the footnote text that provides additional information.}
\end{Verbatim}
\end{code}
Footnotes are limited to one line in print---keep them tight.

\section{Equations}\label{equations}

You can include mathematical equations in your document using \LaTeX{} syntax and use \texttt{\textbackslash{}label\{\}} to reference them later.

\begin{code}{markdown}{}{}
\begin{Verbatim}[commandchars=\\\{\},breaklines, breakanywhere, commandchars=\\\{\}]
\PYZbs{}begin\PYZob{}equation\PYZcb{}
E = mc\PYZca{}2
\PYZbs{}label\PYZob{}eq:einstein\PYZcb{}
\PYZbs{}end\PYZob{}equation\PYZcb{}

As shown in Equation \PYZdl{}\PYZbs{}eqref\PYZob{}eq:einstein\PYZcb{}\PYZdl{}, energy is equal to mass times the speed of light squared.
\end{Verbatim}
\end{code}
For consistency, TeXSmith provides the shorthand \texttt{@[label]} to reference it.

\begin{code}{markdown}{}{}
\begin{Verbatim}[commandchars=\\\{\},breaklines, breakanywhere, commandchars=\\\{\}]
As shown in Equation @[eq:einstein], energy is equal to mass times the speed of light squared.
\end{Verbatim}
\end{code}
\section{Figures}\label{figures}

Give a figure a caption block with an \texttt{id} and reference it anywhere with the
\texttt{@[label]} shorthand. The \texttt{id} is emitted verbatim as the \LaTeX{} \texttt{\textbackslash{}label}.

\begin{code}{markdown}{}{}
\begin{Verbatim}[commandchars=\\\{\},breaklines, breakanywhere, commandchars=\\\{\}]
![\PY{n+nt}{Sample Figure}](\PY{n+na}{image\PYZhy{}url.jpg})

/// caption
    attrs: \PYZob{}id: sample\PYZhy{}figure\PYZcb{}
This is the caption for the figure.
///

As shown in Figure @[sample\PYZhy{}figure], the data illustrates...
\end{Verbatim}
\end{code}
\begin{callout}[callout warning]{Warning}
The \texttt{id} may not contain a colon: \texttt{pymdown-\allowbreak{}extensions} rejects
identifiers such as \texttt{fig:sample-\allowbreak{}figure}, and the whole \texttt{/// caption}
block then silently falls back to plain text (TeXSmith emits a warning
when this happens). Use plain identifiers like \texttt{sample-\allowbreak{}figure}.
\end{callout}
Both web and print outputs number figures automatically, though the actual numbers may differ because each layout floats content differently.

\section{Tables}\label{tables}

Tables present structured data; give them a caption line with a label so you
can reference them later. The \texttt{Table:} line goes directly above the table.

\begin{code}{markdown}{}{}
\begin{Verbatim}[commandchars=\\\{\},breaklines, breakanywhere, commandchars=\\\{\}]
Table: A sample table for cross\PYZhy{}references. \PYZob{}\PYZsh{}tbl:sample\PYZhy{}table\PYZcb{}

| Header 1 | Header 2 |
|\PYZhy{}\PYZhy{}\PYZhy{}\PYZhy{}\PYZhy{}\PYZhy{}\PYZhy{}\PYZhy{}\PYZhy{}\PYZhy{}|\PYZhy{}\PYZhy{}\PYZhy{}\PYZhy{}\PYZhy{}\PYZhy{}\PYZhy{}\PYZhy{}\PYZhy{}\PYZhy{}|
| Cell 1   | Cell 2   |

Check Table @[tbl:sample\PYZhy{}table] for more details.
\end{Verbatim}
\end{code}
\section{Code Block References}\label{code-block-references}

Labelled, cross-referenceable code listings are not implemented yet: fenced
code blocks accept a \texttt{title} (rendered as the listing header) but carry no
label, so there is no \texttt{@[...]} target for them. Track the feature before
relying on it.

\section{Tags and Index Entries}\label{tags-and-index-entries}

Add tags or index entries to associate keywords with specific sections or topics in the document.

\begin{code}{markdown}{}{}
\begin{Verbatim}[commandchars=\\\{\},breaklines, breakanywhere, commandchars=\\\{\}]
This section covers advanced sorting algorithms. \PYZob{}index\PYZcb{}[algorithm]
\end{Verbatim}
\end{code}
See the section [Index / Tags][index-tags] for more details on how to manage index entries.

\section{Naming Conventions}\label{naming-conventions}

Before hypertext, references revolved around numbers: pages, figures, tables, equations.

\begin{code}{text}{}{}
\begin{Verbatim}[commandchars=\\\{\},breaklines, breakanywhere, commandchars=\\\{\}]
1. Section

  Some Text

  Table 7: An example table
         +\PYZhy{}\PYZhy{}\PYZhy{}\PYZhy{}\PYZhy{}\PYZhy{}\PYZhy{}+
         | Table |
         +\PYZhy{}\PYZhy{}\PYZhy{}\PYZhy{}\PYZhy{}\PYZhy{}\PYZhy{}+

         +\PYZhy{}\PYZhy{}\PYZhy{}\PYZhy{}\PYZhy{}\PYZhy{}\PYZhy{}\PYZhy{}+
         | Figure |
         +\PYZhy{}\PYZhy{}\PYZhy{}\PYZhy{}\PYZhy{}\PYZhy{}\PYZhy{}\PYZhy{}+
  Figure 42: An example figure

2. Another Section

  See Table 7 for more details. The Figure 42 illustrates the concept.
  Everything is explained in Section 1.
\end{Verbatim}
\end{code}
\subsection{French}\label{french}

In French, the reference type stays lowercase unless it begins the sentence.

\begin{displayquote}

Voir le tableau 7 pour plus de détails. La figure 42 illustre le concept.
Le tout est expliqué à la section 1.

\end{displayquote}
\subsection{English}\label{english}

In English we capitalize the reference type and skip the article (``Table 7,'' not ``the Table 7'').

\begin{displayquote}

See Table 7 for more details. Figure 42 illustrates the concept.
Everything is explained in Section 1.

\end{displayquote}
\subsection{German}\label{german}

German capitalizes the reference type too.

\begin{displayquote}

Siehe Tabelle 7 für weitere Details. Abbildung 42 veranschaulicht das Konzept.
Alles wird in Abschnitt 1 erklärt.

\end{displayquote}

\end{document}
